% Bilagan till övningen: bara metodkapitlet.  Övningen gör inga egna
% litteratursökningar; varje källa den citerar är belagd i en av veckans
% föreläsningar, och de påståenden som ännu inte har någon källa står i
% avsnittet Fortsatt arbete.  Ingen \chapterprecis här: den hör till
% kapitel som redovisar ett sökresultat, och det här kapitlet gör inte
% det.

\chapter{Att kontrollera ett påstående}
\label{app:verifiera}

Ett lösningsförslag är fullt av påståenden.  De flesta kan du pröva
själv genom att köra programmet --- det är därför varje körning i den
här texten är en riktig körning av just den kod som visas med den.  Men
några påståenden går inte att pröva på det sättet:
\begin{enumerate}
  \item att en fråga om medlemskap kostar i genomsnitt lika mycket
    oavsett hur många element en mängd eller en uppslagslista har, medan
    samma fråga i en lista kostar mer ju längre listan är
    (\cref{sec:epost,sec:banken});
  \item att \mintinline{python}{collections.deque} är byggd för att
    lägga till och ta bort i båda ändarna, till skillnad från en lista
    (\cref{sec:cowsay});
  \item att en nyckel i en uppslagslista måste vara oföränderlig
    (\cref{sec:brak});
  \item att samma startvärde till slumpgeneratorn ger samma följd varje
    gång (\cref{sec:humor});
  \item att en fullständig kontroll av en e-postadress följer en standard
    som tillåter mer än kontrollen i \cref{sec:epost};
  \item att det är bättre att välja den enklaste behållare som klarar de
    operationer programmet behöver (\cref{sec:banken,sec:cowsay});
  \item de didaktiska påståenden som anteckningarna i marginalen gör om
    hur människor lär sig.
\end{enumerate}
Hur vet man om sådant stämmer?  Samma fråga möter dig varje gång du
läser en lärobok, en webbsida eller ett svar från en språkmodell --- och
varje gång du själv ska skriva en rapport.  Den här bilagan beskriver en
arbetsgång för att kontrollera ett påstående, och
\cref{tab:ovning-pastaenden} visar var vart och ett av de uppräknade
påståendena är belagt.  Lägg särskilt märke till det första: en siffra om
hur snabb en operation är ser exakt ut som ett faktum, men gäller bara
under förutsättningar som måste anges.

\begin{table}[!htb]
  \centering
  \caption{Övningens påståenden, som frågor, och var svaret finns.  De
    fyra första är frågor om vad Python gör och besvaras av Pythons egen
    dokumentation; de två därpå är belagda i andra föreläsningar och
    citeras därifrån; den sista raden gäller marginalanteckningarna.}
  \label{tab:ovning-pastaenden}
  \begin{tabular}{@{}p{0.46\linewidth}p{0.44\linewidth}@{}}
    \toprule
    Fråga & Var det är belagt \\
    \midrule
    Kostar en fråga om medlemskap lika mycket oavsett antalet element i
    en mängd eller en uppslagslista, till skillnad från i en lista?
      & Primärkälla: Pythons egen sammanställning av operationernas
        kostnad, läst i fulltext 2026-09-03 för föreläsningen
        \emph{Behållare: Mängder, stackar och köer}.  Ja --- i
        genomsnitt, i CPython; i värsta fall nej. \\
    Är \mintinline{python}{deque} byggd för båda ändarna, och en lista
    inte?
      & Primärkälla: dokumentationen för modulen
        \mintinline{python}{collections}, läst i fulltext 2026-09-03 för
        samma föreläsning. \\
    Måste en nyckel vara oföränderlig?
      & Primärkälla: Pythons handledning, kapitlet \emph{Data
        Structures}, läst i fulltext 2026-09-03 för föreläsningen
        \emph{Behållare: Uppslagslistor}. \\
    Ger samma startvärde samma följd av slumptal?
      & Primärkälla: dokumentationen för modulen
        \mintinline{python}{random}, läst i fulltext 2026-09-03 för
        föreläsningen \emph{Behållare: Ett gissningsspel}. \\
    Tillåter standarden för e-postadresser mer än kontrollen här?
      & Inte belagt.  Se \cref{sec:ovning-fortsatt}. \\
    Ska lösningen hållas enkel (\foreignlanguage{english}{KISS})?
      & Redan belagt i föreläsningen \emph{Funktioner}, bilaga~C: att
        enkelhet är ett etablerat designvärde, medan den vanliga
        upphovsuppgiften inte går att bekräfta. \\
    Måste kontrast upplevas samtidigt för att en aspekt ska kunna
    urskiljas, och förstör en fråga som pekar ut det kritiska den
    urskiljning den ber om?
      & Primärkälla: Marton, \emph{Necessary Conditions of Learning},
        s.~45--47 respektive s.~13 och 91. \\
    \bottomrule
  \end{tabular}
\end{table}

\section{Gör påståendet till en fråga}

Ett påstående går att kontrollera först när det är tydligt vad som skulle
kunna visa att det är \emph{fel}.  Formulera det därför som en fråga med
ett svar: \enquote{kostar det mer att söka i en lista ju längre listan
är?} kan besvaras med ja eller nej, medan \enquote{uppslagslistor är
bra} inte kan det --- ingen tänkbar källa skulle visa att det är fel,
eftersom \enquote{bra} är ett omdöme och inte ett faktum.  Ofta döljer
sig flera påståenden i ett: \enquote{det är bättre att lagra kontona i en
uppslagslista} är minst tre frågor --- om hastighet, om vilka fel
behållaren gör omöjliga, och om vad man förlorar --- och de kräver olika
slags källor.

\section{Välj rätt sorts källa}

Olika frågor besvaras av olika källor.
\begin{description}
  \item[Uppslagsverk] för vad ett ord \emph{betyder}.
  \item[Primärkällan] för var en idé \emph{kommer ifrån}, och för hur ett
    system \emph{beter sig}.  Frågan om vad en av Pythons behållare gör
    besvaras av Pythons egen dokumentation, inte av en lärobok som
    återger den --- och inte av en körning på din egen dator, som visar
    vad just den datorn gjorde just då.
  \item[Översikter] (\foreignlanguage{english}{reviews}) för om något är
    \emph{etablerat}, eller för vad forskningen \emph{sammantaget}
    säger.
  \item[Enskilda studier] för ett specifikt resultat --- med förbehållet
    att en enskild studie bara visar att något observerats en gång.
\end{description}
Var man hittar dem: universitetsbibliotekets söktjänst, och databaser som
Scopus, Web of Science, IEEE Xplore, ACM Digital Library, DBLP (datalogi)
och OpenAlex (öppen, och trots namnet inte detsamma som öppen tillgång).
Google Scholar täcker brett men rankar ogenomskinligt.  För tekniska
standarder --- som den för e-postadresser --- är primärkällan
standarddokumentet självt, och de är fritt tillgängliga.

\section{Sök så att du kan ha fel}

Den som söker på namn hon redan känner till hittar det hon väntar sig.
Sök därför \emph{ämnesdrivet} --- på begrepp, inte författare.  Några
praktiska regler: håll frågorna korta, ställ samma frågor i flera
databaser, sök begreppets alla namn, och sök också efter
\emph{invändningar} genom att lägga till ord som \enquote{critique},
\enquote{limitations} eller \enquote{failure}.  Ett påstående som
överlevt en motsökning är värt mer än ett som aldrig prövats, och en
invändning som hittas blir ett \emph{förbehåll} i svaret, inte ett
nederlag.

\section{Läs i fulltext, och skriv ner hur}

En träff i en databas är inte ett belägg, och sammanfattningen
(\foreignlanguage{english}{abstract}) räcker inte heller: den säger vad
författarna vill ha sagt, inte vad de visat.  Läs hela texten och leta
upp \emph{det ställe} som styrker påståendet --- ett ordagrant citat med
sidnummer.  Räkna aldrig en källa du inte läst.

Skriv sedan ner hur du gick till väga, så att någon annan kan följa
spåret.  I det här materialet står den anteckningen som en kommentar
ovanför varje post i litteraturlistans källfil, med fasta fält:
\begin{description}
  \item[\texttt{CLAIM}] vilket påstående källan ska styrka;
  \item[\texttt{FOUND-VIA}] hur den hittades;
  \item[\texttt{PICKED}] varför just den, framför andra träffar;
  \item[\texttt{QUOTE}] det ordagranna citatet, med sida eller avsnitt;
  \item[\texttt{VERIFIED}] att fulltexten lästs, och var;
  \item[\texttt{COUNTER}] vad motsökningen gav;
  \item[\texttt{DATE}] när.
\end{description}
Det tar ett par minuter per källa, och det är det enda som gör det
möjligt att om ett år svara på frågan \enquote{hur vet du det?}

\section{Dra slutsatsen --- med förbehåll}
\label{sec:ovning-slutsats}

Svaret skrivs ut tillsammans med det som talar emot.  Ett påstående om
hur snabb en operation är gäller till exempel en viss implementation och
ett genomsnittsfall, inte alla tänkbara fall --- och det förbehållet hör
till svaret.  Förbehållen är inte svagheter i svaret.  De \emph{är}
svaret, och de avgör hur påståendet får användas i materialet.  Det är
också därför lösningsförslagen i den här övningen säger när en skillnad
är omätbar i praktiken: ett riktigt påstående som används i fel skala
blir ett felaktigt råd.

\section{En anmärkning om språkmodeller}

En språkmodell kan användas för att \emph{hitta} och \emph{sortera}
källor.  Lägg märke till arbetsfördelningen: varje källa som citeras ska
läsas och väljas av en människa.  Använd gärna språkmodeller för att
hitta och sortera; låt dem aldrig vara det sista ledet.  De kan göra
fel, men det är du själv som måste stå till svars för innehållet.

\section{Fortsatt arbete}
\label{sec:ovning-fortsatt}

Den här bilagan är bara metodkapitlet: övningen har inte gjort någon egen
litteratursökning, eftersom varje källa den citerar redan är läst och
redovisad i en av föreläsningarna.  Fem påståenden i övningen
saknar ändå ett belägg, och de är utpekade med kommentarer i källfilen
(\texttt{contents.nw}) så att de går att hitta:
\begin{enumerate}
  \item Att standarden för e-postadresser tillåter mer än kontrollen i
    \cref{sec:epost} --- bland annat citerade lokaldelar och domäner
    utan punkt.  Primärkällan är RFC 5322, som inte är läst i fulltext.
    Frågan är värd att läsa på plats, eftersom standarden är fritt
    tillgänglig och eftersom de flesta program i världen medvetet
    avviker från den.
  \item Att \mintinline{python}{collections.Counter} räknar förekomster
    och svarar \(0\) på ett värde som aldrig förekommit
    (\cref{sec:bokstaver}).  Dokumentationen för modulen är läst i
    fulltext, men de citat som är nedtecknade täcker bara
    \mintinline{python}{deque}; \mintinline{python}{Counter} behöver ett
    eget citat.
  \item Att Python tar bort \mintinline{python}{__hash__} från en klass
    som definierar \mintinline{python}{__eq__} (\cref{sec:brak}).
    Körningen i avsnittet visar ett fall, men regeln som sådan står i
    språkreferensen, som saknar post i litteraturlistan.
  \item Att \mintinline{python}{functools.total_ordering} härleder de
    övriga jämförelseoperatorerna ur \mintinline{python}{__eq__} och
    \mintinline{python}{__lt__} (\cref{sec:brak}).  Påståendet är en
    hänvisning till dokumentationen och är inte kontrollerat.
  \item Att \mintinline{python}{random.choice} drar ur samma generator
    som \mintinline{python}{random.seed} styr (\cref{sec:humor}).
    Körningen visar att det stämmer, men citaten som är nedtecknade för
    modulen täcker \mintinline{python}{seed} och
    \mintinline{python}{randint}, inte \mintinline{python}{choice}.
\end{enumerate}
Gemensamt för fyra av de fem är att körningen i materialet visar
\emph{ett} fall, medan påståendet gäller en regel.  Det är den vanligaste
formen av obelagt påstående i programmeringsmaterial, och den är värd att
känna igen: att något fungerade en gång är inte detsamma som att det är
garanterat.
